\subsection{Structural Syntax and Function Object Creation}

A Python function is not only a named block of source code. At runtime, it is an object. The \texttt{def} statement creates that object, connects it to a code object, attaches metadata such as its name and documentation string, and then binds a name to it in the current namespace.

This follows the same object-reference model already used throughout the previous chapters: a name is not a C-style storage box. A name points to an object. In this section, the object happens to be executable.

The important distinction is this:

\begin{itemize}
    \item executing a \texttt{def} statement creates a function object;
    \item calling the function later creates a new execution frame;
    \item the body of the function does not run merely because the function was defined.
\end{itemize}

The next subsection will study return behavior and frame termination. Here we focus only on function creation, function objects, code objects, and lambda expressions.

\subsubsection{The \texttt{def} Keyword: Execution-Time Function Object Creation and Name Binding}

The \texttt{def} statement is executed by the interpreter like other statements. When execution reaches it, Python creates a function object and binds the chosen name to that object.

The body is not executed at definition time.

\begin{verbatim}
print("before def")

def say_fact():
    """Print one short fact."""
    print("Chuck Norris counted to infinity. Twice.")

print("after def")
print()

print(f"say_fact object: {say_fact}")
print(f"type(say_fact): {type(say_fact)}")

print()
print("now calling say_fact()")
result = say_fact()

print()
print(f"returned value: {result!r}")
print(f"type(result):   {type(result)}")
\end{verbatim}

Possible output:

\begin{verbatim}
before def
after def

say_fact object: <function say_fact at 0x...>
type(say_fact): <class 'function'>

now calling say_fact()
Chuck Norris counted to infinity. Twice.

returned value: None
type(result):   <class 'NoneType'>
\end{verbatim}

The line \texttt{def say\_fact():} creates the function object. The line \texttt{say\_fact()} calls it.

This means that a function definition is already a runtime event. Python does not merely remember a textual block for later. It constructs an executable object and places a reference to that object under a name.

A function object can be inspected like other objects.

\begin{verbatim}
def say_fact():
    """Print one short fact."""
    print("Chuck Norris counted to infinity. Twice.")

selected_names = [
    "__call__",
    "__name__",
    "__doc__",
    "__defaults__",
    "__kwdefaults__",
    "__annotations__",
    "__code__",
]

print(f"type(say_fact): {type(say_fact)}")
print()

for name in selected_names:
    print(f"{name:<18} {name in dir(say_fact)}")
\end{verbatim}

Possible output:

\begin{verbatim}
type(say_fact): <class 'function'>

__call__          True
__name__          True
__doc__           True
__defaults__      True
__kwdefaults__    True
__annotations__   True
__code__          True
\end{verbatim}

The presence of \texttt{\_\_call\_\_} is what makes the object callable. The programmer writes \texttt{say\_fact()}, but that syntax is backed by the object's call behavior.

The name created by \texttt{def} is still only a name. It can be rebound. The old function object does not become invalid just because another name now points somewhere else.

\begin{verbatim}
def line():
    return "Nobody expects the Spanish Inquisition!"

old_line = line

def line():
    return "It is only a model."

print(f"old_line(): {old_line()}")
print(f"line():     {line()}")

print()
print(f"old_line object: {old_line}")
print(f"line object:     {line}")
print(f"same object:     {old_line is line}")
\end{verbatim}

Possible output:

\begin{verbatim}
old_line(): Nobody expects the Spanish Inquisition!
line():     It is only a model.

old_line object: <function line at 0x...>
line object:     <function line at 0x...>
same object:     False
\end{verbatim}

The second \texttt{def line():} creates a new function object and binds the name \texttt{line} to that new object. The name \texttt{old\_line} still points to the first function object.

So the real structure is:

\begin{verbatim}
first def line():
    creates function object A
    binds name line -> A

old_line = line:
    binds name old_line -> A

second def line():
    creates function object B
    rebinds name line -> B
\end{verbatim}

This is the same naming model already seen for integers, strings, lists, tuples, sets, and dictionaries. The difference is that this object can be executed.

\subsubsection{Function Objects as Runtime Values: \texttt{\_\_name\_\_}, \texttt{\_\_doc\_\_}, \texttt{\_\_defaults\_\_}, \texttt{\_\_kwdefaults\_\_}, and \texttt{\_\_annotations\_\_}}

A function object carries metadata. Some of that metadata comes directly from the source code.

\begin{verbatim}
def build_line(name: str, score: int = 42, *, mood: str = "calm") -> str:
    """Build one formatted display line."""
    return f"{name:<10} score={score:04d} mood={mood}"

print(f"function object: {build_line}")
print(f"type:            {type(build_line)}")
print()

print(f"__name__:        {build_line.__name__!r}")
print(f"__doc__:         {build_line.__doc__!r}")
print(f"__defaults__:    {build_line.__defaults__!r}")
print(f"__kwdefaults__:  {build_line.__kwdefaults__!r}")
print(f"__annotations__: {build_line.__annotations__!r}")

print()
print(build_line("Homer"))
print(build_line("Arthur", 7, mood="confused"))
\end{verbatim}

Possible output:

\begin{verbatim}
function object: <function build_line at 0x...>
type:            <class 'function'>

__name__:        'build_line'
__doc__:         'Build one formatted display line.'
__defaults__:    (42,)
__kwdefaults__:  {'mood': 'calm'}
__annotations__: {'name': <class 'str'>, 'score': <class 'int'>, 'mood': <class 'str'>, 'return': <class 'str'>}

Homer      score=0042 mood=calm
Arthur     score=0007 mood=confused
\end{verbatim}

The metadata can be read directly:

\begin{itemize}
    \item \texttt{\_\_name\_\_} stores the function's internal name;
    \item \texttt{\_\_doc\_\_} stores the docstring, or \texttt{None} if there is no docstring;
    \item \texttt{\_\_defaults\_\_} stores positional default values as a tuple;
    \item \texttt{\_\_kwdefaults\_\_} stores keyword-only default values as a dictionary, or \texttt{None};
    \item \texttt{\_\_annotations\_\_} stores annotations as a dictionary.
\end{itemize}

The default value \texttt{42} is not stored as text. It is stored as an object inside the function object's metadata. This matters later when we discuss default argument evaluation and the mutable default argument trap.

The annotations are also metadata. They do not automatically force runtime type checking.

\begin{verbatim}
def show_value(x: int) -> str:
    return f"value={x!r}"

print(f"annotations: {show_value.__annotations__}")
print(show_value(42))
print(show_value("No soup for you!"))
\end{verbatim}

Possible output:

\begin{verbatim}
annotations: {'x': <class 'int'>, 'return': <class 'str'>}
value=42
value='No soup for you!'
\end{verbatim}

The annotation says that \texttt{x} is intended to be an \texttt{int}, but Python does not enforce that automatically. Tools may read the annotation. The runtime call itself still proceeds unless the function body performs an operation that fails.

Function objects also have a writable attribute dictionary. This should not be overused, but it proves again that a function is a normal runtime object.

\begin{verbatim}
def quote():
    """Return a small quote."""
    return "D'oh!"

quote.source = "The Simpsons"
quote.level = 42

print(f"quote():       {quote()}")
print(f"quote.__dict__: {quote.__dict__}")
print()

print(f"type(quote.__dict__): {type(quote.__dict__)}")
print(f"'keys' in dir:        {'keys' in dir(quote.__dict__)}")
print(f"'items' in dir:       {'items' in dir(quote.__dict__)}")
\end{verbatim}

Possible output:

\begin{verbatim}
quote():       D'oh!
quote.__dict__: {'source': 'The Simpsons', 'level': 42}

type(quote.__dict__): <class 'dict'>
'keys' in dir:        True
'items' in dir:       True
\end{verbatim}

This does not mean ordinary code should use functions as arbitrary storage bags. It only exposes the object model clearly: a function is an object with metadata, behavior, and references to other objects.

\subsubsection{The Anatomy of Function Code Objects: Inspecting Bytecode Attributes (\texttt{\_\_code\_\_}, \texttt{co\_code}, \texttt{co\_varnames}, \texttt{co\_consts}, \texttt{co\_names})}

A function object contains a reference to a code object. The function object is the callable runtime wrapper. The code object contains compiled execution information: local variable names, constants, bytecode, and several other low-level attributes.

\begin{verbatim}
def make_message(name, number=42):
    text = f"{name} says {number}"
    return text.upper()

code_obj = make_message.__code__

print(f"function type: {type(make_message)}")
print(f"code type:     {type(code_obj)}")
print()

selected_names = [
    "co_code",
    "co_varnames",
    "co_consts",
    "co_names",
    "co_argcount",
    "co_filename",
    "co_firstlineno",
]

for name in selected_names:
    print(f"{name:<16} {name in dir(code_obj)}")
\end{verbatim}

Possible output:

\begin{verbatim}
function type: <class 'function'>
code type:     <class 'code'>

co_code          True
co_varnames      True
co_consts        True
co_names         True
co_argcount      True
co_filename      True
co_firstlineno   True
\end{verbatim}

The function object and the code object are different objects. The function object is callable. The code object is not normally called directly by beginner-level Python code.

The most useful code-object fields for inspection are the following:

\begin{verbatim}
def make_message(name, number=42):
    text = f"{name} says {number}"
    return text.upper()

code_obj = make_message.__code__

print(f"co_argcount:    {code_obj.co_argcount}")
print(f"co_varnames:    {code_obj.co_varnames}")
print(f"co_consts:      {code_obj.co_consts}")
print(f"co_names:       {code_obj.co_names}")
print(f"co_firstlineno: {code_obj.co_firstlineno}")
print()

print(f"co_code type:   {type(code_obj.co_code)}")
print(f"co_code length: {len(code_obj.co_code)}")
print(f"co_code hex:    {code_obj.co_code.hex(' ')}")
\end{verbatim}

Possible output:

\begin{verbatim}
co_argcount:    2
co_varnames:    ('name', 'number', 'text')
co_consts:      (None, ' says ')
co_names:       ('upper',)
co_firstlineno: 1

co_code type:   <class 'bytes'>
co_code length: 58
co_code hex:    97 00 7c 00 9b 00 64 01 7c 01 9b 00 9d 03 7d 02 ...
\end{verbatim}

The exact byte sequence is version-dependent. CPython bytecode changes across Python releases. The stable idea is not the exact numbers, but the separation:

\begin{verbatim}
function object:
    callable object with metadata and defaults

code object:
    compiled execution description

co_code:
    raw bytecode bytes

co_varnames:
    local variable names known by the compiled function

co_consts:
    constants embedded in the compiled function

co_names:
    external names referenced by the compiled function
\end{verbatim}

Because \texttt{co\_code} is a \texttt{bytes} object, it can also be inspected as an object.

\begin{verbatim}
def make_message(name, number=42):
    text = f"{name} says {number}"
    return text.upper()

code_bytes = make_message.__code__.co_code

print(f"type(code_bytes): {type(code_bytes)}")
print()

selected_names = ["hex", "__len__", "__getitem__", "count", "index"]

for name in selected_names:
    print(f"{name:<14} {name in dir(code_bytes)}")

print()
print(f"first 8 byte values: {list(code_bytes[:8])}")
print(f"hex view:            {code_bytes[:8].hex(' ')}")
\end{verbatim}

Possible output:

\begin{verbatim}
type(code_bytes): <class 'bytes'>

hex            True
__len__        True
__getitem__    True
count          True
index          True

first 8 byte values: [151, 0, 124, 0, 155, 0, 100, 1]
hex view:            97 00 7c 00 9b 00 64 01
\end{verbatim}

The \texttt{bytes} object is not the source code. It is the low-level instruction stream used by CPython's virtual machine. Reading it directly is not pleasant. The standard \texttt{dis} module displays it in a more readable form.

\begin{verbatim}
import dis

def make_message(name, number=42):
    text = f"{name} says {number}"
    return text.upper()

dis.dis(make_message)
\end{verbatim}

Possible output on a modern CPython version:

\begin{verbatim}
  3           0 RESUME                   0

  4           2 LOAD_FAST                0 (name)
              4 FORMAT_VALUE             0
              6 LOAD_CONST               1 (' says ')
              8 LOAD_FAST                1 (number)
             10 FORMAT_VALUE             0
             12 BUILD_STRING             3
             14 STORE_FAST               2 (text)

  5          16 LOAD_FAST                2 (text)
             18 LOAD_METHOD              0 (upper)
             40 CALL                     0
             50 RETURN_VALUE
\end{verbatim}

The exact instruction list may differ by Python version. The important structure remains the same: the source-level function body has been compiled into instructions. When the function is called, a frame executes those instructions.

A compact way to see the object chain is:

\begin{verbatim}
def make_message(name, number=42):
    text = f"{name} says {number}"
    return text.upper()

print("name -> function object -> code object -> bytecode bytes")
print()

print(f"name value:       {make_message}")
print(f"function type:    {type(make_message)}")
print(f"code object:      {make_message.__code__}")
print(f"code object type: {type(make_message.__code__)}")
print(f"bytecode type:    {type(make_message.__code__.co_code)}")
\end{verbatim}

Possible output:

\begin{verbatim}
name -> function object -> code object -> bytecode bytes

name value:       <function make_message at 0x...>
function type:    <class 'function'>
code object:      <code object make_message at 0x..., file "...", line 1>
code object type: <class 'code'>
bytecode type:    <class 'bytes'>
\end{verbatim}

This is why a Python function should not be imagined as only a label placed on top of source text. It is a runtime object connected to compiled execution data.

\subsubsection{Lambda Expressions: Expression-Level Construction of Anonymous Function Objects}

A \texttt{lambda} expression also creates a function object. The difference is syntactic. \texttt{def} is a statement. \texttt{lambda} is an expression.

Because \texttt{lambda} is an expression, it can appear where a value is expected.

\begin{verbatim}
add_42 = lambda x: x + 42

print(f"add_42 object: {add_42}")
print(f"type(add_42):  {type(add_42)}")
print(f"add_42(8):     {add_42(8)}")
print()

print(f"__name__:      {add_42.__name__!r}")
print(f"__doc__:       {add_42.__doc__!r}")
print(f"__defaults__:  {add_42.__defaults__!r}")
print(f"__code__:      {add_42.__code__}")
\end{verbatim}

Possible output:

\begin{verbatim}
add_42 object: <function <lambda> at 0x...>
type(add_42):  <class 'function'>
add_42(8):     50

__name__:      '<lambda>'
__doc__:       None
__defaults__:  None
__code__:      <code object <lambda> at 0x..., file "...", line 1>
\end{verbatim}

The name \texttt{add\_42} is not the lambda's internal function name. It is only a variable name bound to the function object. The function object's internal name is \texttt{'<lambda>'}.

A lambda still has a code object.

\begin{verbatim}
add_42 = lambda x: x + 42

code_obj = add_42.__code__

print(f"type(add_42):   {type(add_42)}")
print(f"type(code_obj): {type(code_obj)}")
print()

print(f"co_varnames:    {code_obj.co_varnames}")
print(f"co_consts:      {code_obj.co_consts}")
print(f"co_names:       {code_obj.co_names}")
print(f"co_code type:   {type(code_obj.co_code)}")
print(f"co_code hex:    {code_obj.co_code.hex(' ')}")
\end{verbatim}

Possible output:

\begin{verbatim}
type(add_42):   <class 'function'>
type(code_obj): <class 'code'>

co_varnames:    ('x',)
co_consts:      (None, 42)
co_names:       ()
co_code type:   <class 'bytes'>
co_code hex:    97 00 7c 00 64 01 7a 00 00 00 53 00
\end{verbatim}

A \texttt{lambda} is therefore not a special lightweight mathematical object. It is still a Python function object. It has a type, a code object, bytecode, local variable names, constants, and call behavior.

The difference between \texttt{def} and \texttt{lambda} can be shown directly.

\begin{verbatim}
def add_def(x):
    return x + 42

add_lambda = lambda x: x + 42

items = [
    ("add_def", add_def),
    ("add_lambda", add_lambda),
]

for label, fn in items:
    print(f"{label:<12} type={type(fn)}")
    print(f"{'':<12} __name__={fn.__name__!r}")
    print(f"{'':<12} result={fn(8)}")
    print(f"{'':<12} vars={fn.__code__.co_varnames}")
    print(f"{'':<12} consts={fn.__code__.co_consts}")
    print()
\end{verbatim}

Possible output:

\begin{verbatim}
add_def      type=<class 'function'>
             __name__='add_def'
             result=50
             vars=('x',)
             consts=(None, 42)

add_lambda   type=<class 'function'>
             __name__='<lambda>'
             result=50
             vars=('x',)
             consts=(None, 42)
\end{verbatim}

Both objects are functions. Both can be called. Both have code objects. The lambda form is only more restricted: it contains one expression, not a full statement block.

This is valid:

\begin{verbatim}
small = lambda x: x.strip().upper()
print(small("  ni!  "))
\end{verbatim}

Possible output:

\begin{verbatim}
NI!
\end{verbatim}

This is not valid Python:

\begin{verbatim}
bad = lambda x:
    y = x.strip()
    return y.upper()
\end{verbatim}

For multi-step behavior, use \texttt{def}.

\begin{verbatim}
def clean_text(x):
    y = x.strip()
    return y.upper()

print(clean_text("  ni!  "))
\end{verbatim}

Possible output:

\begin{verbatim}
NI!
\end{verbatim}

The practical rule is simple:

\begin{quote}
Use \texttt{def} when the behavior deserves a named, readable block. Use \texttt{lambda} only when a small expression-level function is clearer than a separate definition.
\end{quote}

At the runtime level, both forms create function objects.